% 本文件由 LaTeX 排版 Agent 根据有效 translation.json 自动生成。
% author=Apocrypha; work_id=0812; title=比拉多之死

\chapter{比拉多之死}

% structure: p0001 source=h1 target=omitted reason=page-title-already-rendered-by-parent

罗马皇帝提庇留凯撒身患重病，得知在耶路撒冷有一位名叫耶稣的医生，能以一句话治愈所有疾病。他不知道犹太人和比拉多已将耶稣处死，便命令他的一位名叫沃鲁西亚努斯的朋友说：\Quote{你尽快渡海前去，告诉我的仆人和朋友比拉多，请他将这位医生送来，好使他恢复我昔日的健康。}沃鲁西亚努斯听了皇帝的命令，立即出发，奉命来到比拉多那里。他向比拉多传达了提庇留凯撒托付的话，说：\Quote{罗马皇帝、你的主上提庇留凯撒听说在这城里有一位医生，仅凭言语就能治愈疾病，恳请你将他送去，为他治疗疾病。}比拉多听后非常害怕，因他知道自己出于嫉妒将耶稣处死。比拉多回答这位使者说：\Quote{此人是个作恶者，诱惑了全体民众；因此城中智者开会，我下令把他钉死了。}使者返回旅店，遇见一位名叫韦罗尼加的妇人，她曾是耶稣的朋友。使者说：\Quote{妇人，这城里曾有一位医生，仅凭言语就能治愈病人，犹太人为何将他处死？}她开始哭泣，说：\Quote{哀哉！我的主，我的天主和我的主，比拉多出于嫉妒交出祂，判了罪，并下令把祂钉死。}使者极其忧伤，说：\Quote{我深为悲伤，因为我无法完成我主派我来做的事。}韦罗尼加对他说：\Quote{当我的主四处传道时，我因不情愿地失去祂的同在，便想请人为祂画像，以便在失去祂同在时，祂画像的形象至少能给我安慰。当我带着画布去找画师时，我的主遇见我，问我去哪里。我向祂说明缘由，祂向我要了那块布，然后还给我，上面印着祂可敬面容的像。因此，如果你的主上虔诚地瞻仰祂的面容，必会立即获得健康的恩惠。}使者对她说：\Quote{这样一幅画像能用金银买到吗？}她对他说：\Quote{不能，只能凭虔诚的敬爱之心。因此，我将与你同行，带着画像去给凯撒看，然后再回来。}\par

于是沃鲁西亚努斯与韦罗尼加来到罗马，对皇帝提庇留说：\Quote{你所渴望的耶稣，已被比拉多和犹太人处以不义的死刑，出于嫉妒钉在十字架上。现有一位贵妇与我同来，带来了耶稣本人的画像；你若虔诚地注视它，必立即恢复健康。}凯撒便命人用丝绸铺路，将画像呈到他面前；他一看见，便恢复了昔日的健康。\par

因此，奉凯撒之命，般雀·比拉多被逮捕并押解到罗马。凯撒听说比拉多抵达罗马，对他怒火中烧，命人将他带来。比拉多随身带着耶稣的无缝长衣，在皇帝面前穿着它。皇帝一见他，便放下一切怒气，立即起身迎接他，无法对他说任何严厉的话。那个在缺席时显得如此可畏和暴怒的人，如今在眼前却变得温和。皇帝打发他走后，立刻又对他大发雷霆，痛斥他是个该死的人，说他活在世上是不光彩的。皇帝立即又命人召回他，起誓宣称他该死。但一见到比拉多，又立即向他致意，抛开心中的所有暴戾。众人都感到惊异，皇帝自己也惊奇：为何在比拉多不在时对他如此愤怒，而在场时却无法对他说一句严词。于是，因着天主的感动，或某个基督徒的建议，皇帝命人脱去那件长衣，立即又恢复了对他的暴怒。皇帝对此大为惊讶，有人告诉他那件长衣属于主耶稣。于是皇帝下令将他囚禁，等到智者会议商讨如何处置他。几日后，对比拉多作出判决，应处以最可耻的死刑。比拉多听后，用自己的刀自尽，以这样一死结束了生命。\par

凯撒得知比拉多的死讯，说：\Quote{他确实死于最可耻的死亡，他的手竟没有放过自己。}于是将他绑在一块大石头上，沉入台伯河。但邪恶污秽的灵在他邪恶污秽的身体里一同欢跃，在水中搅动，并在空中可怕地制造闪电、暴风、雷霆和冰雹，使众人陷入极度的恐惧。因此，罗马人将他从台伯河捞出，嘲弄地拖到维也纳，沉入罗讷河。因为维也纳被称为\OriginalTerm{地狱之路}{Via Gehennæ}，因那时它是个被咒诅的地方。但那里又有邪灵作祟，行同样的奇事。那些人无法忍受这样鬼魔的侵扰，便将这咒诅的器皿从自己那里移开，送到洛桑尼亚境内埋葬。那里的人看到自己被上述的侵扰所困扰，又将他移走，沉入一个群山环绕的深坑中。据某些人所说，直到今日，那里仍有魔鬼的诡计在翻腾涌出。\par

\SourceRecord{Translated by Alexander Walker.；Ante-Nicene Fathers；Vol. 8.；Edited by Alexander Roberts, James Donaldson, and A. Cleveland Coxe.；Buffalo, NY: Christian Literature Publishing Co.,；1886.；<http://www.newadvent.org/fathers/0812.htm>.}
